\cvsection{Grants and Contracts}

\cvsubsection{Active}

\cvrow{2026 - Present}{\textbf{1UG3CA309489-01 (PI Davis)} \newline Role: Co-Investigator / Biostatistician \newline The BUNDLE study will adapt proven single-cancer outreach strategies into a multicomponent intervention targeting all USPSTF-recommended screenings (colorectal, breast, cervical, and lung), engaging 18+ primary care clinics and 1,600+ patients across Oregon and North Carolina. Using a six-year cluster randomized trial design, the study aims to achieve a 10\% or greater increase in screening completion while evaluating cost-effectiveness and the multilevel factors that influence screening behavior.}
\cvrow{2026 - Present}{\textbf{3OT2OD038368-01S1 (PI Davis)} \newline Role: Biostatistician \newline The PRaCTICe Network Research Hub partners with rural clinics across six Western states to support primary care research, education, and technical assistance. Through four cores, Community Engagement, Recruitment \& Trials, Data Coordination, and Communications \& Dissemination, it facilitates participant recruitment, codesigns studies, and strengthens rural research infrastructure. The network aims to enhance primary care engagement, implement research advances, and improve health outcomes in rural communities.}
\cvrow{2025 - Present}{\textbf{5R01AG082931-02 (PI Dorr)} \newline Role: Biostatistician \newline The MC COMPARE project uses aggregated EHR data via the eCarePlan application to improve care for patients with multiple chronic conditions. It aims to quantify risks, support personalized decisions, and enhance clinical trial outcomes by harmonizing data from multiple sources and engaging patients and providers.}

\cvsubsection{Completed}

\cvrow{2024 - 2025}{\textbf{1R25GM150166-01 (PI Marriott)} \newline Role: Biostatistician \newline This study explores factors associated with biomedical workforce professional development. It measures the development of STEM beliefs and psychosocial factors over time, including how courses, training experiences, and/or instructional practices influence decision-making and career trajectories. This study explores the representation of professional self in portfolio materials for diverse trainees and professionals.}
